\documentclass[../main.tex]{subfiles}
\begin{document}

\textit{Kế thừa kiến trúc đã mô tả ở Chương 3, chương này phân tích cơ sở nhận thức và lý thuyết đằng sau từng quyết định thiết kế chính của CogMem. Mục tiêu của chương không phải là chứng minh mỗi thành phần luôn thắng trên mọi benchmark, mà là chỉ ra vì sao các thành phần đó là những lựa chọn biểu diễn và truy hồi có cơ sở khoa học.}

\section{Phân tích lý thuyết}

\subsection{Cơ sở khoa học nhận thức cho Memory Graph (6 Networks)}

\subsubsection{Semantic Memory và World Network}

Tulving (1972) định nghĩa semantic memory là hệ thống lưu trữ kiến thức chung về thế giới, tương đối tách rời khỏi bối cảnh trải nghiệm cụ thể. Điều này là cơ sở cho \textbf{World Network} trong CogMem: những fact như nghề nghiệp, kỹ năng, quan hệ khách quan hay tri thức miền có thể được lưu dưới dạng world fact mà không cần gắn chặt với một episode riêng lẻ.

Về mặt hệ thống, việc tách world ra khỏi experience giúp giảm nhập nhằng giữa ``kiến thức về người dùng'' và ``một lần người dùng kể lại điều gì đó''. Đây là khác biệt quan trọng cho cả retrieval lẫn synthesis.

\begin{figure}[H]
\centering
\textbf{[FIGURE PLACEHOLDER]}\\
Hình 4.1: Semantic Memory và World Network.
\caption{Semantic Memory và World Network}
\label{fig:ch4_semantic_memory}
\end{figure}

\subsubsection{Episodic Memory và Experience Network}

Episodic memory theo Tulving (1983) lưu trữ các sự kiện cá nhân gắn với thời gian và hoàn cảnh. Đây là nền tảng lý thuyết cho \textbf{Experience Network}. Khi hệ thống cần trả lời các câu hỏi ``đã xảy ra khi nào'', ``trong lần nào'', hoặc ``ở giai đoạn nào của câu chuyện'', các experience node đóng vai trò như mốc sự kiện cơ bản của trí nhớ hội thoại.

Việc duy trì các trường thời gian như \textit{occurred\_start} và \textit{occurred\_end} không chỉ là chi tiết kỹ thuật; nó phản ánh trực tiếp ý tưởng rằng episodic recall cần một trục thời gian để tổ chức các sự kiện đã xảy ra.

\subsubsection{Habit Memory và Habit Network}

Habit memory thường được liên hệ với các mẫu hành vi lặp lại, ít phụ thuộc vào deliberation tức thời. Trong CogMem, \textbf{Habit Network} đóng vai trò như nơi biểu diễn các regularities kiểu ``user thường làm gì'', ``hay chọn gì'', ``phản ứng quen thuộc trong ngữ cảnh nào''.

Từ góc độ nhận thức, habit khác với event đơn lẻ ở chỗ nó tổng quát hóa qua nhiều lần lặp. Từ góc độ hệ thống, điều này giải thích vì sao habit không nên bị xem như một experience bình thường. Một experience nói về một lần xảy ra; một habit nói về một pattern bền hơn.

Implementation hiện tại dùng các liên kết S-R ở mức thực dụng để neo habit với các experience liên quan, thay vì mô hình reinforcement learning hoàn chỉnh. Dù vậy, hướng biểu diễn vẫn giữ đúng tinh thần của habit memory: tách pattern lặp lại ra khỏi event đơn lẻ.

\begin{figure}[H]
\centering
\textbf{[FIGURE PLACEHOLDER]}\\
Hình 4.2: Habit Memory và Habit Network.
\caption{Habit Memory và Habit Network}
\label{fig:ch4_habit_memory}
\end{figure}

\subsubsection{Prospective Memory và Intention Network}

Prospective memory được mô tả như khả năng ``nhớ để thực hiện một việc trong tương lai''. Đây là nền tảng cho \textbf{Intention Network}. Trong hội thoại dài hạn, các phát biểu như ``tôi định học Rust'', ``tuần tới tôi sẽ nộp hồ sơ'', hay ``tôi chưa làm việc đó'' không nên bị ép hết vào experience hoặc world, vì chúng chứa một loại trạng thái khác: kế hoạch chưa hoàn tất.

CogMem mô hình hóa ý tưởng này bằng \textit{intention\_status}. Một intention có thể ở trạng thái \textit{planning}, \textit{fulfilled} hoặc \textit{abandoned}. Quan trọng hơn, transition không được hiểu đơn giản là ``mọi trạng thái đều là edge''. Trong implementation hiện tại:
\begin{itemize}
    \item \texttt{fulfilled\_by} nối intention với experience khi kế hoạch được thực hiện,
    \item \texttt{triggered} nối experience sang intention khi một sự kiện khơi ra kế hoạch mới,
    \item \texttt{enabled\_by} nối giữa các intention,
    \item còn \textit{abandoned} được giữ như một trạng thái của node, không phải một edge type riêng.
\end{itemize}

\begin{table}[H]
\centering
\caption{Intention lifecycle và typed transition semantics}
\label{tab:intention_lifecycle}
\begin{tabular}{|p{3.0cm}|p{4.0cm}|p{4.4cm}|}
\hline
\textbf{Tình huống} & \textbf{Biểu diễn trên intention node} & \textbf{Transition semantics} \\
\hline
Lập kế hoạch mới & status = planning & chưa cần transition \\
\hline
Kế hoạch được thực hiện & status có thể cập nhật; experience mới được tạo & fulfilled\_by: intention $\rightarrow$ experience \\
\hline
Sự kiện gợi ra kế hoạch mới & intention planning mới xuất hiện & triggered: experience $\rightarrow$ intention \\
\hline
Kế hoạch phụ thuộc kế hoạch khác & hai intention cùng tồn tại & enabled\_by: intention $\rightarrow$ intention \\
\hline
Kế hoạch bị hủy bỏ & status = abandoned & không cần edge riêng; trạng thái nằm trên node \\
\hline
\end{tabular}
\end{table}

Các ablation nhắm đích thực hiện sau này trong dự án cho thấy intention network không phải lúc nào cũng tạo lợi thế đồng đều. Tuy nhiên, điều đó không phủ định cơ sở của nó; ngược lại, nó gợi ý rằng prospective representation đặc biệt quan trọng ở các trường hợp mà nội dung ``plan-not-done'' không được world hoặc experience hấp thụ tự nhiên.

\begin{figure}[H]
\centering
\textbf{[FIGURE PLACEHOLDER]}\\
Hình 4.3: Intention lifecycle và prospective memory discharge.
\caption{Intention Lifecycle}
\label{fig:ch4_intention_lifecycle}
\end{figure}

\subsubsection{Theory of Event Coding (TEC) và Action-Effect Network}

Theory of Event Coding (TEC) nhấn mạnh rằng hành động và hệ quả cảm nhận được của hành động là hai mặt của cùng một event representation. Đây là nền tảng cho \textbf{Action-Effect Network} của CogMem. Khi một hệ thống muốn trả lời câu hỏi causal như ``vì sao cách làm đó hiệu quả?'' hoặc ``hành động nào đã dẫn tới kết quả này?'', việc có một node chuyên biệt mang bộ ba \textit{precondition -- action -- outcome} là hữu ích hơn so với chỉ lưu một câu world fact chung chung.

Điểm lý thuyết quan trọng ở đây là causal rule không chỉ là kiến thức về thế giới; nó còn là kiến thức về \textit{can thiệp}. Nói ``X là đúng'' khác với nói ``nếu ở trạng thái P, làm A sẽ dẫn tới O''. Action-effect network giữ nguyên cấu trúc can thiệp đó.

Các thí nghiệm nhắm đích ở workload agentic về sau cho thấy giá trị của action-effect cũng mang tính điều kiện: nó đặc biệt hữu ích khi causal rule không bị tái mã hóa dễ dàng thành một world fact kiểu ``X được giải quyết bằng Y''. Điều này củng cố cách nhìn của chương này: action-effect là một lớp biểu diễn cần thiết cho một số loại causal knowledge, chứ không phải luôn là lớp biểu diễn duy nhất có thể hoạt động.

\begin{figure}[H]
\centering
\textbf{[FIGURE PLACEHOLDER]}\\
Hình 4.4: TEC và Action-Effect Network.
\caption{Theory of Event Coding và Action-Effect Network}
\label{fig:ch4_tec_network}
\end{figure}

\subsection{Cơ sở khoa học nhận thức cho Raw Snippet}

Fuzzy Trace Theory của Brainerd và Reyna cho rằng con người lưu đồng thời ít nhất hai mức dấu vết trí nhớ:
\begin{itemize}
    \item \textbf{gist}: phần tóm tắt ngữ nghĩa, đủ để nhận diện bản chất của sự kiện;
    \item \textbf{verbatim}: phần chi tiết nguyên văn, cần thiết khi phải nhớ chính xác con số, trích dẫn hoặc thứ tự cụ thể.
\end{itemize}

CogMem phản ánh ý tưởng này bằng \textbf{fact text} và \textbf{raw\_snippet}. Fact text hỗ trợ retrieval vì ngắn gọn và dễ nhúng; raw\_snippet hỗ trợ synthesis vì giữ lại chi tiết nguyên bản. Ở implementation hiện tại, snippet có thể được inline trực tiếp vào generation prompt khi cần tăng grounding cho các câu hỏi đòi hỏi độ chính xác cao. Ngay cả khi policy sinh không luôn hiển thị snippet, việc bảo toàn nó vẫn là một yêu cầu kiến trúc.

\begin{figure}[H]
\centering
\textbf{[FIGURE PLACEHOLDER]}\\
Hình 4.5: Fuzzy Trace Theory và 2-layer schema.
\caption{Fuzzy Trace Theory và Raw Snippet}
\label{fig:ch4_fuzzy_trace}
\end{figure}

\subsection{Cơ sở khoa học nhận thức cho SUM + Cycle Guards}

\subsubsection{Episodic Buffer của Baddeley}

Baddeley mô tả episodic buffer như không gian tích hợp tạm thời, nơi nhiều nguồn thông tin từ các hệ con bộ nhớ khác nhau có thể được kết hợp thành một biểu diễn làm việc thống nhất. Ý tưởng này gần với cách CogMem dùng SUM spreading activation: không chọn duy nhất một đường mạnh nhất, mà tích lũy nhiều dấu vết bằng chứng có liên quan.

Từ góc nhìn này, SUM không chỉ là một mẹo kỹ thuật để tăng recall. Nó là một lựa chọn mô hình hóa hợp lý khi coi truy hồi như một quá trình xây dựng lại episode từ nhiều mảnh tín hiệu rải rác.

\begin{figure}[H]
\centering
\textbf{[FIGURE PLACEHOLDER]}\\
Hình 4.6: Episodic Buffer của Baddeley.
\caption{Episodic Buffer của Baddeley}
\label{fig:ch4_episodic_buffer}
\end{figure}

\subsubsection{Hodgkin-Huxley và Local Refractory Period}

Trong thần kinh học, refractory period ngăn neuron vừa phát xung bị kích hoạt lại ngay lập tức. Ý tưởng này được chuyển hóa trong CogMem thành guard chống ping-pong giữa hai node hoặc một vòng lặp cục bộ:

\begin{equation}
\mathrm{refractory}(u) =
\begin{cases}
0, & \text{nếu } u \text{ vừa được kích hoạt} \\
1, & \text{ngược lại}
\end{cases}
\end{equation}

Về ý nghĩa tính toán, guard này biến propagation thành một quá trình tiến về phía trước thay vì dao động quanh cùng một cụm node.

\subsubsection{Wilson-Cowan và Saturation Threshold}

Các mô hình động học thần kinh như Wilson-Cowan cho thấy hoạt động mạng cần những cơ chế phi tuyến để tránh bùng nổ năng lượng và giữ hệ ổn định. CogMem phản ánh điều này qua \textit{saturation threshold}: activation được chặn trên thay vì tăng vô hạn qua mỗi lần cộng dồn.

Điều này quan trọng ở retrieval vì mục tiêu không chỉ là ``đẩy một node lên càng cao càng tốt'', mà còn là giữ được một tập candidate đa dạng để tầng fusion và synthesis còn có thể khai thác.

\subsection{Cơ sở khoa học nhận thức cho Adaptive Query Routing}

Adaptive Query Routing của CogMem có thể được xem như một dạng \textit{attentional selection} ở cấp hệ thống. Tùy theo bản chất truy vấn, hệ thống chuyển trọng tâm sang các nguồn bằng chứng khác nhau:

\begin{itemize}
    \item \textbf{Temporal queries}: tăng trọng số cho temporal channel và temporal constraint.
    \item \textbf{Causal queries}: tăng trọng số cho graph evidence và action-effect signals.
    \item \textbf{Prospective queries}: ưu tiên planning intention và transition evidence.
    \item \textbf{Preference queries}: ưu tiên evidence liên quan đến habit, opinion và hành vi lặp lại.
    \item \textbf{Multi-hop queries}: tăng vai trò của graph traversal để nối nhiều mảnh bằng chứng.
\end{itemize}

Về mặt nhận thức, đây là phép mô phỏng ý tưởng rằng bộ não không dùng cùng một chiến lược tìm kiếm cho mọi câu hỏi. Về mặt kỹ thuật, nó giải thích vì sao CogMem cần query analyzer, trọng số dung hợp theo loại truy vấn, và các bước hậu xử lý như prospective filtering hoặc evidence-priority boosting.

\begin{figure}[H]
\centering
\textbf{[FIGURE PLACEHOLDER]}\\
Hình 4.7: Attentional selection và Adaptive Query Routing.
\caption{Attentional Selection và Adaptive Routing}
\label{fig:ch4_attentional_selection}
\end{figure}


\vspace{1em}
\noindent\textit{Sau khi cơ sở lý thuyết của các quyết định thiết kế đã được làm rõ, Chương 5 sẽ chuyển sang phần đánh giá thực nghiệm để xem những giả thuyết này biểu hiện như thế nào trên dữ liệu benchmark.}

\end{document}
